%==============================================================================
% Artigo AP2 - IBM3119 (Sistemas Embarcados)
% Modelo: SBrT (baseado em IEEEtran, que e a base do template do SBrT).
%
% >>> IMPORTANTE <<<
% O enunciado pede o MODELO DO SBrT2026. No Overleaf:
%   1) Crie um projeto a partir do template oficial do SBrT2026 (procure
%      "SBrT" na galeria de templates do Overleaf, ou use o .zip do simposio).
%   2) Cole o conteudo das secoes abaixo (\section ... ) no main.tex do template.
% Este arquivo usa IEEEtran para JA COMPILAR no Overleaf mesmo sem o .sty do
% SBrT, servindo de rascunho completo. Troque a classe pelo template oficial
% antes da entrega final.
%
% Limite: 3 paginas. Idioma: ingles (incentivo de +0,5 do enunciado).
% Lacunas a preencher com SEUS numeros estao marcadas com  %% TODO
%==============================================================================
\documentclass[conference]{IEEEtran}
\usepackage[utf8]{inputenc}
\usepackage{graphicx}
\usepackage{amsmath}
\usepackage{url}
\usepackage{booktabs}

\begin{document}

\title{An IoT-Based Eye-Tracking Framework for Driver Attention Monitoring and
Traffic Accident Mitigation}

\author{
\IEEEauthorblockN{Yago Duarte \and %% TODO: demais autores do grupo
}
\IEEEauthorblockA{Ibmec --- IBM3119 Embedded Systems\\
%% TODO: e-mails
}
}

\maketitle

\begin{abstract}
Driver drowsiness is a leading cause of traffic accidents, especially among
long-haul truck drivers. This work presents an IoT-based eye-tracking framework
that monitors driver attention in real time from a single camera mounted on the
interior rear-view mirror. The system uses the MediaPipe Face Landmarker to
extract 478 facial landmarks and combines a geometric Eye Aspect Ratio (EAR)
signal with the PERCLOS fatigue metric and an automatic per-driver calibration
step. A continuous drowsiness score drives a three-level escalating audible
alarm, and a lightweight web dashboard streams the live state for remote
monitoring. We quantitatively compare two eye-closure detection methods --- a
geometric EAR threshold and the neural eye-blink blendshape produced by
MediaPipe --- reaching an F1-score of %% TODO and %% TODO respectively, while
running at %% TODO\,FPS on a commodity CPU. Results show the framework is a
viable, low-cost driver-monitoring solution for embedded deployment.
\end{abstract}

\begin{IEEEkeywords}
Driver drowsiness, eye tracking, EAR, PERCLOS, MediaPipe, IoT, embedded vision.
\end{IEEEkeywords}

%==============================================================================
\section{Introduction}
%==============================================================================
Fatigue and drowsiness at the wheel are responsible for a significant share of
fatal road crashes %% TODO: cite a statistic (e.g., national traffic agency).
The risk is higher for professional truck drivers, who face long shifts and
monotonous routes. Driver Monitoring Systems (DMS) based on computer vision can
detect early signs of drowsiness --- prolonged eye closure, slow blinks,
yawning --- and warn the driver in time to prevent an accident.

This paper describes an embedded, IoT-oriented eye-tracking framework that
adapts a study-oriented gaze-tracking system into a drowsiness-detection
system for vehicles. The camera is relocated from a frontal desktop position to
the interior rear-view mirror, as in commercial DMS. The main contributions are:
(i) a real-time pipeline combining EAR, PERCLOS and automatic calibration;
(ii) an escalating multi-level alarm proportional to a continuous drowsiness
score; (iii) a web dashboard for live monitoring; and (iv) a quantitative
comparison between a geometric and a learning-based eye-closure detector.
The source code is available at~\cite{repo}.

%==============================================================================
\section{System Description and Problem}
%==============================================================================
The system targets \emph{driver drowsiness detection}. Unlike the original
study system, which estimated \emph{where} the user was looking, here the goal
is to estimate \emph{how closed} the eyes are over time and to react before the
driver falls asleep. The camera is mounted on the rear-view mirror, giving an
off-axis view of the face (Fig.~\ref{fig:setup}). The system must run in real
time on a CPU, work for different drivers, and tolerate the mirror angle and
varying illumination inside the cabin.

\begin{figure}[t]
\centering
%% TODO: inserir foto do posicionamento da camera (retrovisor) e/ou print do dashboard
%\includegraphics[width=\columnwidth]{figures/setup.png}
\caption{Camera placement on the interior rear-view mirror and the live
monitoring dashboard.}
\label{fig:setup}
\end{figure}

%==============================================================================
\section{Proposed Method}
%==============================================================================
\textbf{Landmark extraction.} Each frame is processed by the MediaPipe Face
Landmarker (Tasks API), which returns 478 normalized facial landmarks and a set
of blendshape scores, including \texttt{eyeBlinkLeft} and \texttt{eyeBlinkRight}.

\textbf{Method A --- Geometric EAR.} For each eye we compute the Eye Aspect
Ratio from six landmarks:
\begin{equation}
\mathrm{EAR}=\frac{\lVert p_2-p_6\rVert+\lVert p_3-p_5\rVert}{2\,\lVert p_1-p_4\rVert}.
\end{equation}
A 3-second \emph{automatic calibration} measures the open-eye EAR of the driver,
so all thresholds become \emph{relative} ($\mathrm{ratio}=\mathrm{EAR}/\mathrm{EAR}_\text{open}$).
This makes the method robust to the mirror angle and to per-person differences.

\textbf{Method B --- Learning-based blink.} As an alternative detector we use the
MediaPipe eye-blink blendshape directly, declaring the eye closed when the blink
score exceeds a threshold.

\textbf{Temporal fusion and scoring.} We track PERCLOS (percentage of eye
closure over a 60\,s window) and a continuous drowsiness score in $[0,100]$ that
\emph{increases} while the eyes are closed/half-closed and \emph{decreases} when
they reopen, plus a penalty for yawns (detected via the Mouth Aspect Ratio).
The score maps to three alert levels (attention/alert/danger) that drive an
\emph{escalating} audible alarm (higher pitch and shorter interval as the score
grows).

\textbf{IoT dashboard.} A Flask server streams the annotated video (MJPEG) and
the metrics (JSON polling) to a browser dashboard showing the score, level,
EAR, PERCLOS, FPS and a live history chart, enabling remote monitoring.

%==============================================================================
\section{Results}
%==============================================================================
\subsection{Quantitative Results}
We recorded labelled sessions where the driver alternated between alert and
eyes-closed states, logging per-frame EAR, blink score and a manual ground-truth
label. Table~\ref{tab:results} compares the two eye-closure detectors at their
best-F1 operating thresholds; Fig.~\ref{fig:cm} shows the confusion matrices and
Fig.~\ref{fig:bars} the metric comparison.

%% TODO: rode  python evaluate.py  e cole a tabela gerada (figures/metrics_table.tex)
%\begin{table}[t]
\centering
\caption{Comparacao quantitativa dos metodos de deteccao de olho fechado
(1259 quadros rotulados).}
\label{tab:results}
\begin{tabular}{lccccc}
\hline
Metodo & Limiar & Acuracia & Precisao & Recall & F1 \\
\hline
A: EAR (geometrico) & 0.41 & 0.932 & 0.941 & 0.916 & 0.928 \\
B: Blendshape (ML) & 0.53 & 0.931 & 0.933 & 0.923 & 0.928 \\
\hline
\end{tabular}
\end{table}

\begin{table}[t]
\centering
\caption{Quantitative comparison of the two eye-closure detection methods
(REPLACE with values from evaluate.py).}
\label{tab:results}
\begin{tabular}{lccccc}
\toprule
Method & Thr. & Acc. & Prec. & Rec. & F1 \\
\midrule
A: EAR (geometric) & %% TODO & %% TODO & %% TODO & %% TODO & %% TODO \\
B: Blendshape (ML) & %% TODO & %% TODO & %% TODO & %% TODO & %% TODO \\
\bottomrule
\end{tabular}
\end{table}

\begin{figure}[t]
\centering
%\includegraphics[width=\columnwidth]{figures/confusion_matrices.pdf}
\caption{Confusion matrices for Method A (EAR) and Method B (blendshape).}
\label{fig:cm}
\end{figure}

\begin{figure}[t]
\centering
%\includegraphics[width=\columnwidth]{figures/metrics_comparison.pdf}
\caption{Accuracy, precision, recall and F1 for both methods.}
\label{fig:bars}
\end{figure}

We also measured real-time performance: the full pipeline runs at
%% TODO\,FPS on a %% TODO CPU, and the alarm latency (time from eye closure to
audible warning) is below %% TODO\,s.

\subsection{Qualitative Results}
Qualitatively, the automatic calibration allowed the same thresholds to work
for the frontal and the rear-view camera positions. The system remained stable
under %% TODO (e.g., moderate lighting changes / glasses). The escalating alarm
was perceived as %% TODO during the demonstration. Failure cases included
%% TODO (e.g., very low light, extreme head pose, occlusion by sunglasses).

%==============================================================================
\section{Conclusion}
%==============================================================================
We presented a low-cost, IoT-based eye-tracking framework for driver drowsiness
detection suitable for embedded deployment. Combining a calibrated EAR signal,
PERCLOS and an escalating alarm, the system detects sustained eye closure in
real time and was validated by comparing a geometric and a learning-based
detector. Future work includes night-time/IR operation, deployment on a
Raspberry Pi with a GPIO buzzer, head-pose estimation and a larger labelled
dataset with multiple drivers.

%==============================================================================
\begin{thebibliography}{1}
\bibitem{repo} %% TODO: insira a URL do repositorio GitHub (Atividade 1)
Project repository. [Online]. Available: \url{https://github.com/yagoduarte0/...}

\bibitem{fatigue} H. Assump\c{c}\~ao, R. P. Fernandes, A. B. N. Barbosa, and
P. H. P. Gusman, ``Vision-based fatigue detection using a smartphone,'' in
\emph{XLIII Brazilian Symposium on Telecommunications and Signal Processing
(SBrT)}, Natal, RN, Brazil, 2025.

\bibitem{mediapipe} C. Lugaresi \emph{et al.}, ``MediaPipe: A framework for
building perception pipelines,'' \emph{arXiv:1906.08172}, 2019.
\end{thebibliography}

\end{document}
